\title{FlashAttention-3: Fast and Accurate Attention with Asynchrony and Low-precision}

\begin{abstract}
Attention mechanisms are fundamental to the Transformer architecture, yet they represent a major computational bottleneck for large-scale language models and tasks requiring extended context lengths. The \textsc{FlashAttention} approach previously introduced methods to accelerate attention on GPUs by reducing memory access overheads. Nonetheless, prior work such as \textsc{FlashAttention-2} has not fully harnessed the advanced features of recent GPU hardware, with utilization rates reaching only 35\% on H100 GPUs. In this work, we introduce three principal strategies for improving attention efficiency on Hopper GPUs: leveraging the asynchronous capabilities of both Tensor Cores and TMA to (1) facilitate concurrent computation and data transfer through warp specialization, (2) alternate between block-wise matrix multiplication and softmax calculations, and (3) employ block quantization and incoherent computation utilizing hardware-accelerated FP8 low-precision support. Experimental results indicate that our proposed method, \textsc{FlashAttention-3}, provides a 1.5-2.0$\times$ speedup on H100 GPUs, with FP16 performance reaching as high as 740 TFLOPs/s (corresponding to 75\% hardware utilization), and FP8 performance approaching 1.2 PFLOPs/s. Additionally, we show that FP8 \textsc{FlashAttention-3} exhibits numerical errors that are 2.6$\times$ lower than those produced by standard FP8 attention implementations.
\end{abstract}

\section{Introduction}
\label{sec:intro}

Within the Transformer architecture~\citep{vaswani2017attention}, the attention mechanism emerges as the chief source of computational cost, largely because calculating self-attention scores between queries and keys involves operations that scale quadratically with respect to the sequence length. Enhancing the scalability of attention to handle extended contexts stands to facilitate advancements such as more sophisticated modeling and reasoning over several lengthy documents~\citep{guo2021longt5,shaham2022scrolls,peng2023yarn}, navigating extensive code repositories~\citep{roziere2023code, li2023starcoder}, and supporting new modalities—including high-definition images~\citep{chen2022scaling}, audio signals~\citep{gulati2020conformer}, and videos~\citep{ho2022video}. It could also enable innovative use cases, for instance, enabling interaction with prolonged user histories~\citep{sun2019bert4rec} and facilitating agents operating over extended temporal horizons~\citep{yao2022react}. Consequently, there has been growing momentum in efforts to accelerate attention under long-context scenarios, utilizing methods such as approximate algorithms~\citep{katharopoulos2020transformers,choromanski2020rethinking,
tay2020efficient}, software-level improvements (\citep{rabe2021self,dao2022flashattention,kwon2023efficient}), and even the pursuit of fundamentally different architectural frameworks~\citep{peng2023rwkv,sun2023retentive,gu2023mamba}.

This study extends the approach introduced by \citet{dao2022flashattention}, who proposed exact-attention methods that explicitly leverage the GPU’s architectural features and execution patterns in the design of their algorithms. In \citep{dao2022flashattention}, Dao et al. presented \textsc{FlashAttention}, which implemented a new tiling mechanism that accelerates attention computation by combining all related operations within a single GPU kernel, thereby eliminating costly intermediate transfers to and from global memory.

Later, \citet{dao2023flashattention2} advanced the algorithm to \textsc{FlashAttention-2}, restructuring it to enable parallelization over the sequence length and processing blocks of the key and value matrices in the forward pass. This adjustment aimed to better utilize GPU resources by increasing occupancy and balancing computational workload. Despite these enhancements, we note that \textsc{FlashAttention-2} still demonstrates relatively low hardware utilization on newer GPU architectures when compared to highly optimized GEMM kernels—achieving, for instance, only about 35\% utilization versus the 80–90\% achieved by GEMM on the Hopper H100 GPU. This performance gap may stem in part from differences at the implementation level, such as reliance on Ampere-generation instructions instead of those specific to Hopper when operating on Tensor Cores. Recent efforts, including ThunkerKitten~\citep{spector2024thunder} and cuDNN 9~\citep{cudnn9}, illustrate that incorporating Hopper-native instructions and adopting tile-based abstractions can not only accelerate attention computations but also reduce implementation complexity.

At a more fundamental level, the \textsc{FlashAttention-2} algorithm operates within a basic synchronous framework, without leveraging asynchronous execution or low-precision arithmetic in its structure. Asynchronous processing arises from hardware innovations tailored to speed up key ML workload components: certain specialized hardware modules conduct matrix multiplications (such as Tensor Cores) or handle memory transfers (like the Tensor Memory Accelerator, or TMA), operating independently from conventional CUDA cores that handle logic, integer, and floating-point tasks. The adoption of reduced numerical precision—for example, FP8 in Hopper or FP4 in Blackwell, following the earlier use of FP16 (introduced with Pascal in 2017) and BF16 (from Ampere in 2020)—has been validated as a method for achieving two- to fourfold increases in throughput, without extra power or silicon cost. The benefits of such hardware capabilities in Hopper are discussed in \cref{subsec:hardware}. 
The key technical obstacle is to refactor \textsc{FlashAttention-2} to exploit these hardware advancements: employing asynchrony entails overlapping the computation of matmul and softmax, despite their sequential data dependencies, while utilizing low precision demands careful handling to control quantization error—particularly to address feature outliers in LLMs~\citep{dettmers2208llm, sun2024massive}.

With this objective in mind, we introduce \textsc{FlashAttention-3}, which integrates and develops three novel concepts aimed at enhancing performance on the latest GPU architectures:\footnote{Our findings are presented with reference to NVIDIA's Hopper architecture. Nonetheless, the proposed algorithm remains applicable to any GPU platform equipped with advanced asynchronous processing and support for low-precision operations.}

\begin{enumerate}

\item \textbf{Producer-Consumer asynchrony:} We introduce a warp-specialized software pipelining approach that leverages the asynchronous nature of data transfers and Tensor Core computation by assigning data producers and consumers to different warps. This design choice enhances the pipeline’s capacity to mask both memory access and instruction dispatch delays within the algorithm.
\item \textbf{Hiding softmax under asynchronous block-wise GEMMs:} We conceal the relatively slower non-GEMM operations required by softmax—such as floating point multiply-adds and exponentials—by overlapping their execution with asynchronous WGMMA GEMM instructions. To accomplish this, we redesign the \textsc{FlashAttention-2} algorithm so that sequential dependencies between the softmax step and GEMM computations are avoided. For instance, in our 2-stage variant, as softmax processes one scores matrix block, the WGMMA instruction simultaneously and asynchronously computes the subsequent block.
\item \textbf{Hardware-accelerated low-precision GEMM:} We modify the forward pass to support GEMM on FP8 Tensor Cores, resulting in nearly a twofold increase in the observed TFLOPs/s. Adapting to FP8 requires addressing the disparate memory layout requirements for WGMMA, particularly in the arrangement of FP32 accumulators and FP8 input matrices. To minimize the accuracy degradation inherent in lowering to FP8, we apply strategies such as block quantization and incoherent computation.
\end{enumerate}

To assess the empirical effectiveness of our approach, we conduct a series of benchmarks for \textsc{FlashAttention-3} on the H100 SXM5 GPU, exploring various parameter configurations. Our results demonstrate that (1) in the forward pass, FP16 achieves a performance boost of 1.5-2.0$\times$ over \textsc{FlashAttention-2} (with peak throughput up to 740 TFLOPs/s), and in the backward pass achieves 1.5-1.75$\times$ speedup; (2) FP8 mode approaches 1.2 PFLOPs/s; and (3) at long sequence lengths, FP16 surpasses competing methods, while FP8 offers competitive results\footnote{Specifically, FP8 \textsc{FlashAttention-3} leads for head dimension 64, matches for head dimensions 128 and 256 without causal masking, and lags when causal masking is enabled.} compared to the state-of-the-art attention kernel in NVIDIA's cuDNN. Further analysis shows that FP16 \textsc{FlashAttention-3} maintains the same numerical error profile as \textsc{FlashAttention-2} and offers improved accuracy over standard attention due to the use of FP32 for intermediate computations such as softmax rescaling. Additionally, with block quantization and incoherent processing, FP8 \textsc{FlashAttention-3} achieves 2.6$\times$ higher accuracy compared to the standard per-tensor quantized attention, particularly in the presence of outlier features.

We have released \textsc{FlashAttention-3} under a permissive license\footnote{\textsc{FlashAttention-3}
  is available at \texttt{https://github.com/Dao-AILab/flash-attention}}, and we intend to incorporate it into the PyTorch and Hugging Face ecosystems in order to maximize accessibility and support for researchers and practitioners.

\section{Background: Multi-Head Attention and GPU Characteristics}
\label{sec:background}

\subsection{Multi-Head Attention}
\label{subsec:multi_head_attn}

Consider the input sequences for a single attention head, denoted as $\mathbf{Q}, \mathbf{K}, \mathbf{V} \in \mathbb{R}^{N \times d}$, representing queries, keys, and values, respectively, where $N$ stands for the sequence length and $d$ corresponds to the dimensionality of the head. The resulting attention output $\mathbf{O}$ can then be determined by:

\begin{equation*}
  \mathbf{S} = \alpha \mathbf{Q} \mathbf{K}^\top \in \mathbb{R}^{N \times N}, \quad \mathbf{P} = \mathrm{softmax}(\mathbf{S}) \in \mathbb{R}^{N \times N}, \quad \mathbf{O} = \mathbf{P}\mathbf{V} \in \mathbb{R}^{N \times d},
\end{equation*}

Here, $\mathrm{softmax}$ operates across each row, and it is common to use a scaling factor $\alpha = 1/\sqrt{d}$. To mitigate issues with numerical precision caused by exponentiation, we typically subtract $\mathrm{rowmax}(\mathbf{S})$ from $\mathbf{S}$. In the context of multi-head attention (MHA), individual heads possess distinct query, key, and value projection matrices; as a result, these computations are carried out in parallel for all heads and batches, yielding the complete output tensor.

Consider $\phi$ as a scalar-valued loss function, and denote the gradient operator by $\mathbf{d}(-) = \partial \phi / \partial (-)$. Provided with the gradient of the output, $\mathbf{dO} \in \mathbb{R}^{N \times d}$, we utilize the chain rule to derive the gradients $\mathbf{dQ}$, $\mathbf{dK}$, and $\mathbf{dV}$ as shown below:

\begin{align*}
  \mathbf{dV} &= \mathbf{P}^\top \mathbf{dO} \in \mathbb{R}^{N \times d} \\
  \mathbf{dP} &= \mathbf{dO} \mathbf{V}^\top \in \mathbb{R}^{N \times N} \\
  \mathbf{dS} &= \mathrm{dsoftmax} (\mathbf{dP}) \in \mathbb{R}^{N \times N} \\
  \mathbf{dQ} &= \alpha \mathbf{dS} \mathbf{K} \in \mathbb{R}^{N \times d} \\
  \mathbf{dK} &= \alpha \mathbf{dS}^\top \mathbf{Q} \in \mathbb{R}^{N \times d},
\end{align*}

In this context, we observe that $\mathbf{d}s = (\mathrm{diag}(p) - p p^\top)\mathbf{d}p$, where $p = \mathrm{softmax}(s)$ with respect to the vector $s$. The notation $\mathrm{dsoftmax}(\mathbf{dP})$ denotes applying this expression to each row individually. For the backward computation in MHA, this process is similarly parallelizable over all heads and batch instances.

\subsection{GPU hardware characteristics and execution model}
\label{subsec:hardware}

We discuss components of the GPU execution model pertinent to \textsc{FlashAttention-3}, emphasizing the NVIDIA Hopper architecture as a representative example.

\paragraph{Memory hierarchy:} The GPU utilizes a tiered structure for its memory system, in which the amount of storage available decreases as the bandwidth increases (\cref{tab:gpu-hierarchy})\footnote{\citet{luo2024benchmarking} details a shared memory bandwidth of 128 bytes per clock cycle for each SM; to estimate total bandwidth, this is multiplied by 132 SMs and the boost clock frequency of 1830 MHz.}. The global memory (GMEM), or HBM, consists of off-chip DRAM shared among all streaming multiprocessors (SMs). Access to GMEM is buffered by an intermediate on-chip L2 cache, which operates transparently. Beyond this, each SM features a limited-capacity on-chip cache called shared memory (SMEM), which is both highly banked and explicitly managed by the programmer. At the deepest level of the hierarchy, every SM contains its own register file.

\paragraph{Thread hierarchy:} The organizational structure of GPU programming centers on groups of execution entities termed threads. Arranged from smallest to largest, this hierarchy consists of individual threads, followed by warps (comprising 32 threads), warpgroups (made up of 4 sequential warps), threadblocks (also known as cooperative thread arrays or CTAs), threadblock clusters (introduced in Hopper), and finally, grids.

The two hierarchies maintain a strong interconnection. Threads belonging to a single CTA are assigned concurrently to the same SM, while multiple CTAs within one cluster are scheduled together on a single GPC. All threads of a CTA can directly access SMEM, but each individual thread possesses up to 256 registers (RMEM) that remain exclusive to it.

\paragraph{Asynchrony and warp-specialization:}

GPUs function as throughput-oriented processors that depend on both concurrent operations and asynchronous execution to mask delays associated with memory access and computation. To facilitate asynchronous data transfer between GMEM and SMEM, Hopper introduces a specialized hardware component known as the Tensor Memory Accelerator (TMA) \cite[\S7.29]{cuda}. In addition, differentiating itself from earlier designs like Ampere, Hopper’s Tensor Core—utilized through the warpgroup-scoped WGMMA instruction \cite[\S9.7.14]{ptx}—also operates asynchronously and is capable of directly fetching operands from shared memory.

Asynchronous hardware capabilities enable the implementation of warp-specialized kernels, wherein warps within a CTA are assigned exclusively to either producer or consumer tasks—focusing solely on data transfer or computation, respectively. This separation generally enhances the compiler’s potential to optimize instruction scheduling \citep{warp-specialization-2011}. Furthermore, Hopper introduces the ability to dynamically adjust the allocation of registers among warpgroups through \verb|setmaxnreg| \cite[\S9.7.17.1]{ptx}, permitting warps executing MMAs to utilize a greater portion of RMEM compared to those assigned only to TMA operations, which can be effectively handled by a single thread.

\paragraph{Low-precision number formats:}
\label{sec:low-precision-gpu}
Contemporary GPUs incorporate dedicated hardware to enhance performance on low-precision calculations. As an illustration, the WGMMA instruction can utilize Hopper's FP8 Tensor Cores, achieving double the per-SM throughput relative to FP16 or BF16 operations.

Nonetheless, utilizing FP8 WGMMA properly requires careful attention to the required operand memory layouts. Consider the case of a GEMM operation multiplying $A \times B^{\top}$, with $A$ being an $M \times K$ matrix and $B$ being an $N \times K$ matrix. The operand $A$ or $B$ is said to be \emph{mn-major} if elements are stored contiguously along the outer $M$ or $N$ dimension, while it is labeled \emph{k-major} if contiguity exists along the inner $K$ dimension. For FP16 WGMMA, SMEM inputs can be provided in either mn-major or k-major order. In contrast, FP8 WGMMA restricts inputs in SMEM to the k-major layout exclusively. This constraint becomes especially problematic in cases such as attention mechanisms, where one might attempt to fuse multiple consecutive GEMMs in one kernel: the mismatch between the accumulator layout in FP32 and the required layout of the FP8 operands hinders seamless invocation of chained FP8 WGMMA operations.

Regarding attention mechanisms, such layout constraints necessitate adaptations to the FP8 algorithm’s architecture; these adjustments are detailed in \cref{sec:algofp8}.

\subsection{Standard Attention and Flash Attention} In accordance with \citet{dao2022flashattention}, we refer to \textbf{standard attention} as the conventional approach for computing attention on GPUs, in which the intermediate matrices $\mathbf{S}$ and $\mathbf{P}$ are fully written to and read from HBM. The essential innovation in \textsc{FlashAttention} lies in utilizing a localized softmax reduction, thereby eliminating the overhead of storing and accessing these intermediates by consolidating the entire attention operation within a single kernel. This localized softmax is implemented by lines \ref{code-ws:softmax_start}-\ref{code-ws:softmax_end} of the consumer mainloop in \cref{alg:flash3_wgmma_ws_only}, together with rescaling blocks of $\mathbf{O}$. A straightforward proof that this approach correctly computes $\mathbf{O}$ is provided in \cite[\S 2.3.1]{dao2023flashattention2}.

\section{FlashAttention-3: Algorithm}
\label{sec:algo}

This section details the \textsc{FlashAttention-3} algorithm. Our discussion centers on the forward pass, while the backward pass is covered in~\cref{sec:algo_ws_bwd}. We begin by outlining the incorporation of warp-specialization alongside a circular SMEM buffer into the foundational structure of \textsc{FlashAttention-2}. Subsequently, we discuss leveraging WGMMA’s asynchronous capabilities to establish an overlapped, two-stage GEMM-softmax pipeline. Lastly, we present the adjustments required for FP8 support, focusing on both layout compliance and improvements in accuracy achieved through block quantization and incoherent processing.

\subsection{Producer-Consumer asynchrony through warp-specialization and pingpong scheduling}
\label{sec:algo_ws}

\paragraph{Warp-specialization}
Similar to \textsc{FlashAttention-2}, the forward computation in \textsc{FlashAttention-3} exhibits extensive parallelism across batch elements, attention heads, and positions within the query sequence. Accordingly, it is appropriate to describe the algorithm at the CTA level, where each tile $\mathbf{Q}_i$ of the query matrix is processed to generate the respective tile $\mathbf{O}_i$ in the output matrix. For clarity, we first outline the warp-specialization approach using a circular SMEM buffer, omitting any overlap between GEMM and softmax operations for now. Let $d$ represent the head dimension and $N$ the sequence length, with a chosen query block size $B_r$ partitioning $\mathbf{Q}$ into $T_r = \lceil \frac{N}{B_r} \rceil$ blocks labeled $\mathbf{Q}_1, \ldots, \mathbf{Q}_{T_r}$.

\begin{algorithm}[H]
    \caption{\small\label{alg:flash3_wgmma_ws_only}\textsc{FlashAttention-3} forward computation \textbf{excluding} intra-consumer overlap -- CTA perspective}
    \begin{algorithmic}[1]
\REQUIRE Matrices $\mathbf{Q}_i \in \mathbb{R}^{B_r \times d}$ and $\mathbf{K}, \mathbf{V} \in \mathbb{R}^{N \times d}$ stored in HBM, key block size $B_c$, with $T_c = \lceil \frac{N}{B_c} \rceil$.
\STATE Set up a pipeline mechanism to handle barrier synchronization using an $s$-stage circular buffer in shared memory.
\IF {current warpgroup is a producer}
\STATE Free a predefined number of registers.
\STATE Initiate loading of $\mathbf{Q}_i$ from HBM into shared memory.
\STATE Upon finishing the load, signal consumer that $\mathbf{Q}_i$ is available.
\FOR{$0 \le j < T_c$}
    \STATE Await until the $(j\,\%\,s)$th pipeline stage has been processed by consumer.
    \STATE Initiate load operations for $\mathbf{K}_j, \mathbf{V}_j$ from HBM into the shared memory at pipeline stage $(j\,\%\,s)$.
    \STATE Once load is complete, notify consumers that $\mathbf{K}_j, \mathbf{V}_j$ are ready.
\ENDFOR
\ELSE \STATE Allocate a preset number of registers depending on the quantity of consumer warps.
\STATE Locally initialize $\mathbf{O}_i = (0) \in \mathbb{R}^{B_r \times d}$ as well as $\ell_i, m_i = (0), (-\infty) \in \mathbb{R}^{B_r}$.
\STATE Wait until $\mathbf{Q}_i$ appears in shared memory.
\FOR{$0 \le j < T_c$}
\STATE Wait until $\mathbf{K}_j$ is available in shared memory.
\STATE Perform computation $\mathbf{S}_i^{(j)} = \mathbf{Q}_i \mathbf{K}_j^T$ using SS-GEMM. Signal completion and pause. \label{alg:ws_only_gemm1}
\STATE Preserve previous $m_i$ as $m_i^{\mathrm{old}}$ and update $m_i = \mathrm{max}(m_i^{\mathrm{old}}, \mathrm{rowmax}(\mathbf{S}_i^{(j)}))$. \label{code-ws:softmax_start}
\STATE Compute intermediate softmax values: $\widetilde{\mathbf{P}}_i^{(j)} = \mathrm{exp}(\mathbf{S}_i^{(j)} - m_i)$, and update $\ell_i = \mathrm{exp}(m_i^{\mathrm{old}} - m_i) \ell_i + \mathrm{rowsum}(\widetilde{\mathbf{P}}_i^{(j)})$. \label{code-ws:softmax_end}
\STATE Wait for $\mathbf{V}_j$ data in shared memory.
\STATE Update outputs: $\mathbf{O}_i = \mathrm{diag}(\mathrm{exp}(m_i^{\mathrm{old}} - m_i))^{-1} \mathbf{O}_i + \widetilde{\mathbf{P}}_i^{(j)} \mathbf{V}_j$, utilizing RS-GEMM. Signal completion and pause. \label{alg:ws_only_gemm2}
\STATE Unlock the $(j\,\%\,s)$th pipeline stage for the producer.
\ENDFOR
\STATE Finalize outputs: $\mathbf{O}_i = \mathrm{diag}(\ell_i)^{-1} \mathbf{O}_i$ and $L_i = m_i + \log(\ell_i)$.
\STATE Store the resulting $\mathbf{O}_i$ and $L_i$ into HBM as the $i$th segment of $\mathbf{O}$ and $L$.
\ENDIF
\end{algorithmic}
\end{algorithm}

In our implementation of \cref{alg:flash3_wgmma_ws_only} targeting Hopper, we utilize \verb|setmaxnreg| for the purpose of allocating and deallocating registers, employ TMA to transfer $\mathbf{Q}_i$ as well as $\{ \mathbf{K}_j, \mathbf{V}_j \}_{0 \leq j < T_c}$, and make use of WGMMA for executing the GEMM operations within the consumer mainloop. The SS and RS prefixes are used to specify whether shared memory or register file is used as the source of the first operand, respectively. In analyzing the control flow of \cref{alg:flash3_wgmma_ws_only}, it is important to recognize that TMA loads are asynchronous and therefore do not block other loads in progress. Additionally, during the initial $s$ iterations of the producer mainloop, waits are omitted, as the buffer is still in the process of being filled.

\paragraph{Pingpong scheduling} The asynchronous execution of WGMMA and TMA, combined with the use of warp-specialization, creates the possibility for overlapping the softmax calculation in one warpgroup with the GEMM task in a different warpgroup. This approach is compelling given the stark contrast in throughput between matmul and non-matmul instructions on state-of-the-art hardware accelerators. To illustrate, consider the H100 SXM5 GPU, which offers 989 TFLOPS for FP16 matmul operations, yet delivers only 3.9 TFLOPS for special function computations like the exponential\footnote{According to the CUDA programming guide, each streaming multiprocessor (SM) can perform 16 special function operations per clock cycle. To arrive at the 3.9 TFLOPS figure, this is multiplied by 132 SMs and a 1830 MHz clock frequency.} required in softmax. For an FP16 attention forward pass with a head dimension of 128, the workload demands 512 times more matmul floating-point operations than exponentials, but the throughput for exponentials is 256 times less, resulting in exponential computations consuming up to 50\% as much time as matmul operations. This imbalance is exacerbated with FP8: while matmul throughput is doubled, the capacity for exponentials does not increase, further amplifying the throughput disparity.

Given that the exponential operation is handled by a dedicated hardware component (the multi-function unit), it is preferable for the computation of exponentials to overlap with the Tensor Cores' execution of matrix multiplications. To achieve this, we introduce synchronization barriers (\texttt{bar.sync} instructions), ensuring that the GEMM operations (specifically, GEMM1—$\mathbf{P} \mathbf{V}$ from one iteration, and GEMM0—$\mathbf{Q} \mathbf{K}^\top$ from the subsequent iteration) executed by warpgroup 1 are initiated prior to those of warpgroup 2. This scheduling ensures that, while warpgroup 2 performs its GEMMs, the softmax for warpgroup 1 can be computed concurrently. Subsequently, the responsibilities alternate: warpgroup 2 switches to executing softmax as warpgroup 1 resumes its GEMMs, resulting in a ``pingpong'' scheduling mechanism. This pattern is visualized in~\cref{fig:pingpong_scheduling}. Although actual scheduling behavior does not always align perfectly with the idealized depiction in the figure, empirical results indicate a tangible performance gain; for example, FP16 forward pass throughput increases from approximately 570 TFLOPS to 620–640 TFLOPS for head dimension 128 and sequence length 8192.

\paragraph{Attention variants} For both multi-query attention~\citep{shazeer2019fast} and grouped query attention~\citep{ainslie2023gqa}, our method aligns with that of \textsc{FlashAttention-2}, modifying tensor indexing so as to prevent redundant copies of $\mathbf{K}$ and $\mathbf{V}$ within HBM.

\subsection{Intra-warpgroup overlapping GEMMs and softmax}

It is possible to interleave certain operations from the softmax computation with operations from the GEMMs, even when restricted to a single warpgroup. Here, we present a method for achieving this overlap.

Within the attention algorithm, the core loop presents sequential dependencies among its operations, which hinder the potential for intra-iteration parallel execution. Specifically, the (local) softmax computation (spanning lines \ref{code-ws:softmax_start} to \ref{code-ws:softmax_end}) depends on receiving $\mathbf{S}_i^{(j)}$ from the initial GEMM operation, and the second GEMM requires the output $\widetilde{\mathbf{P}}_i^{(j)}$ from softmax. Consequently, the wait statements at lines \ref{alg:ws_only_gemm1} and \ref{alg:ws_only_gemm2} in \cref{alg:flash3_wgmma_ws_only} enforce a strict order between softmax and GEMM executions. Nonetheless, these execution barriers can be alleviated by employing a pipeline structure across loop iterations, enabled through supplementary register buffers. Building on this strategy, we introduce a two-stage\footnote{It is important to emphasize that while the maximum number of pipeline stages is limited by the number $s$ of circular SMEM buffer stages, they do not have to coincide.} pipelining method that overlaps GEMM and softmax computations:

\begin{algorithm}[H]
    \caption{\small\label{alg:flash3_wgmma}\textsc{FlashAttention-3} consumer warpgroup forward pass}
    \begin{algorithmic}[1]
      \REQUIRE  Matrices $\mathbf{Q}_i \in \mathbb{R}^{B_r \times d}$ and $\mathbf{K}, \mathbf{V} \in \mathbb{R}^{N \times d}$ are present in HBM, with the key block size specified as $B_c$ and $T_c = \lceil \frac{N}{B_c} \rceil$.
      \STATE Allocate the predetermined set of registers according to the number of consumer warps.
      \STATE Initialize the on-chip variables as follows: set $\mathbf{O}_i = (0) \in \mathbb{R}^{B_r \times d}$ and $\ell_i, m_i = (0), (-\infty) \in \mathbb{R}^{B_r}$.
      \STATE Ensure that $\mathbf{Q}_i$ and $\mathbf{K}_0$ are available in shared memory before proceeding.
      \STATE Calculate $\mathbf{S}_{\mathrm{cur}} = \mathbf{Q}_i \mathbf{K}_0^T$ by utilizing WGMMA. Commit the operation and wait for completion.
      \STATE Free the $0$th buffer stage used for $\mathbf{K}$.
      \STATE Derive $m_{i}$, $\tilde{\mathbf{P}}_{\mathrm{cur}}$, and $\ell_{i}$ from $\mathbf{S}_{\mathrm{cur}}$, and rescale $\mathbf{O}_i$ as needed.
      \FOR{$1 \le j < T_c - 1$}
        \STATE Await the loading of $\mathbf{K}_j$ into shared memory.
        \label{code:mainloop_start}
        \STATE Using WGMMA, calculate $\mathbf{S}_{\mathrm{next}} = \mathbf{Q}_i \mathbf{K}_{j}^T$. Commit the computation but continue without waiting. \label{code:first_wgmma}
        \STATE Hold until $\mathbf{V}_{j-1}$ is fully loaded in shared memory.
        \STATE Update $\mathbf{O}_{i} = \mathbf{O}_{i} + \tilde{\mathbf{P}}_{\mathrm{cur}} \mathbf{V}_{j-1}$ through WGMMA, committing immediately but not waiting. \label{code:second_wgmma}
        \STATE Wait for the completion of the WGMMA operation for $\mathbf{Q}_i \mathbf{K}_{j}^T$.
        \STATE From $\mathbf{S}_{\mathrm{next}}$, compute $m_{i}$, $\tilde{\mathbf{P}}_{\mathrm{next}}$, and $\ell_{i}$. \label{code:softmax}
        \STATE Await the completion of the WGMMA calculation for $\tilde{\mathbf{P}}_{\mathrm{cur}} \mathbf{V}_{j-1}$, then perform rescaling on $\mathbf{O}_i$.
        \STATE Release the $(j\,\%\,s)$th buffer stage for $\mathbf{K}$, as well as the $(j-1\,\%\,s)$th buffer stage for $\mathbf{V}$.
        \STATE Assign $\mathbf{S}_{\mathrm{next}}$ to $\mathbf{S}_{\mathrm{cur}}$ for the next iteration.
        \label{code:mainloop_end}
      \ENDFOR \STATE Wait until $\mathbf{V}_{T_c - 1}$ has been loaded into shared memory.
      \STATE Using WGMMA, finalize $\mathbf{O}_{i} = \mathbf{O}_{i} + \tilde{\mathbf{P}}_{\mathrm{last}} \mathbf{V}_{T_c - 1}$, committing and waiting for the result.

\STATE Epilogue: Adjust $\mathbf{O}_{i}$ by applying the scaling factor $m_{i}$. Determine $L_{i}$ utilizing both $m_{i}$ and $\ell_{i}$.
Store $\mathbf{O}_{i}$ and $L_{i}$ in HBM, assigning them as the $i$-th segments of $\mathbf{O}$ and $L$, respectively.
\end{algorithmic}
\end{algorithm}

\cref{alg:flash3_wgmma} substitutes for the consumer component of \cref{alg:flash3_wgmma_ws_only} to form the full \textsc{FlashAttention-3} algorithm when operating at FP16 precision. Conceptually, WGMMA serves here as shorthand for asynchronous GEMM execution. Inside the mainloop (from line \ref{code:mainloop_start} through line \ref{code:mainloop_end}), the second WGMMA invocation in iteration $j$ (line \ref{code:second_wgmma}) executes concurrently with the softmax steps of the subsequent iteration $j+1$ (line \ref{code:softmax}).

Although the pipelined approach depicted previously presents potential performance improvements in theory, a number of real-world considerations must be taken into account:

\paragraph{Compiler reordering} The pseudocode suggests an idealized sequence of operations, yet in practice, the compiler (NVCC) frequently reorganizes instructions to enhance performance. Such reordering can interfere with the intended interleaving of WGMMA and non-WGMMA operations, potentially causing behavior that deviates from the intended pipeline and diminishing the anticipated benefits. Examination of the SASS output, however, confirms that the generated code successfully achieves the expected instruction overlap (see Section~\ref{sec:2-stage-sass}).

\paragraph{Register pressure} Achieving maximum efficiency depends on limiting register spilling. Nonetheless, the 2-stage pipeline demands extra registers to hold temporary results and to preserve state across stages. In particular, it is necessary to retain an additional $\mathbf{S}_{\mathrm{next}}$ in registers, increasing register usage by $B_r \times B_c \times \text{sizeof}(\text{float})$ for each threadblock. This added requirement may compete with other optimizations such as using larger block sizes, which themselves place significant pressure on register usage. Consequently, it is important to assess these competing demands using empirical profiling to identify the best balance.

\paragraph{3-stage pipelining} Building upon the previously discussed 2-stage pipeline, we introduce a 3-stage version aimed at overlapping the second WGMMA with the softmax computation. This extension could further improve utilization of Tensor Cores but comes at the cost of needing even more registers for the additional pipeline stage, thereby making the compromise between increasing the tile size and deepening the pipeline more complex. Comprehensive details and performance analysis for this 3-stage approach are provided in ~\cref{sec:3-stage}.

\subsection{Low-precision with FP8}
\label{sec:algofp8}

\textbf{Efficiency: layout transformations.} Executing the forward pass of \textsc{FlashAttention-3} with FP8 precision introduces extra complications related to layout compatibility that are not present when using FP16.

To begin, it is important to highlight that the input tensors $\mathbf{Q}$, $\mathbf{K}$, and $\mathbf{V}$ are usually stored with contiguity along the head dimension. However, meeting the k-major requirement imposed by FP8 WGMMA for the second GEMM operation necessitates that $\mathbf{V}$—specifically, the tiles of $\mathbf{V}$ brought into SMEM—be contiguous with respect to the sequence length dimension. Since the TMA load is unable to alter the axis of contiguity, we are faced with two choices: (1) perform a transpose of $\mathbf{V}$ in GMEM as a pre-processing stage, or (2) conduct an in-kernel transposition of the tiles of $\mathbf{V}$ subsequent to loading them into SMEM. For the first approach, there are two further sub-options: (1a) the transpose can be integrated into the epilogue of an earlier stage, such as rotary embedding, or (1b) a dedicated pre-processing transpose kernel\footnote{A high-performance transpose kernel can reach bandwidths approaching the device’s limits~\citep{colfax_cutlass_transpose_2024}.} can be used to swap the strides between the sequence length and head dimensions. Nonetheless, (1a) poses significant challenges for standard library integration, while (1b) incurs excessive memory traffic, making it impractical for memory-constrained scenarios like inference.

In the case of FP8 \textsc{FlashAttention-3}, we choose to implement option (2). For the in-kernel transpose operation, we leverage the LDSM (\verb|ldmatrix|) and STSM (\verb|stmatrix|) instructions, which enable a group of threads within a warp to cooperatively load data from SMEM into RMEM and store from RMEM to SMEM, operating on blocks of 128 bytes at a time.\footnote{According to the PTX documentation, LDSM and STSM are intended for copying $8 \times 8$ matrices with 16-bit elements \cite[\S 9.7.13.4.15-16]{ptx}; by packing two 8-bit values together, these instructions can also be utilized in the context of FP8. However, since the transposed variants of LDSM/STSM do not unpack 8-bit pairs, some register-level data movement must be interposed between LDSM and STSM for true tile transposition. Details are omitted here.} These instructions offer efficient register utilization, making it practical to run them within the producer warpgroup, and they also allow for transposed memory transfers. Additionally, once the initial iteration completes, we can schedule the transpose operation for the subsequent $\mathbf{V}$ tile to overlap with the execution of the two WGMMAs, which process the previous $\mathbf{V}$ and the current $\mathbf{K}$ tile.

Second, we note that, in contrast to FP16, the organization of the FP32 accumulator in an FP8 WGMMA does not align with the register layout expected for operand A. Fragments of these distinct layouts are illustrated in \cref{fig:rmem_accum} and \cref{fig:rmem_operand}, which display the sequence in which register entries are allocated per thread. Utilizing byte permutation instructions enables us to rearrange the output from the first WGMMA’s accumulator so that it conforms to the format required by the subsequent WGMMA, aligning with the structure of the $\mathbf{V}$ tile generated by the in-kernel transpose. Concretely, and referring to \cref{fig:rmem_accum}, we reorder the sequence to
$$\{ \verb|d0 d1 d4 d5 d2 d3 d6 d7| \},$$
with this register rearrangement being repeated every 8 bytes. From the perspective of the logical shape of the $\mathbf{P}$ tile, this operation effectively permutes its columns (for instance, columns $0189$ are now relocated to the leading four columns). To ensure that WGMMA subsequently produces the desired output tile, we can arrange for the in-kernel transpose to apply a compatible row permutation to the $\mathbf{V}$ tile.\footnote{This flexibility made possible by executing the transpose within the kernel removes the need to use shuffle instructions to reassign register values across threads, as previously detailed in~\citep{colfax_fp8_flashattention_2024}.}

\textbf{Accuracy: block quantization and incoherent processing.} In the FP8 (e4m3) representation, only 3 bits are allocated to the mantissa, and 4 bits are designated for the exponent. This allocation leads to greater numerical inaccuracies when compared to formats such as FP16 or BF16. Additionally, large-scale models often contain outlier values~\citep{dettmers2208llm, sun2024massive} that significantly exceed the magnitude of most elements, which complicates the quantization process. A common strategy is to utilize per-tensor scaling~\citep{micikevicius2022fp8}, assigning a single scaling factor to each tensor (for example, one each for $\mathbf{Q}$, $\mathbf{K}$, and $\mathbf{V}$). To alleviate the increased numerical errors in FP8 attention, we adopt two mitigation strategies:

\begin{enumerate}

\item \textbf{Block quantization}: In this approach, a single scalar is assigned to represent each block; thus, for every one of $\mathbf{Q}$, $\mathbf{K}$, and $\mathbf{V}$, the tensor is divided into blocks of dimensions $B_r \times d$ or $B_c \times d$, and each block undergoes quantization independently. This process can be efficiently integrated with operations immediately preceding the attention calculation—such as rotary embedding—without introducing latency, given that rotary embedding is constrained by memory bandwidth. As \textsc{FlashAttention-3} inherently works with block structures, it is possible to rescale each segment of $\mathbf{S}$ accordingly to compensate for block quantization effects, all without incurring additional computational overhead.

\item \textbf{Incoherent processing}: To mitigate the influence of outliers, we pre-multiply both $\mathbf{Q}$ and $\mathbf{K}$ by a random orthogonal matrix $\mathbf{M}$ prior to performing FP8 quantization. The orthogonality condition, $\mathbf{M} \mathbf{M}^\top = I$, ensures that $(\mathbf{Q} \mathbf{M}) (\mathbf{K} \mathbf{M})^\top = \mathbf{Q} \mathbf{K}^\top$—meaning this transformation leaves the output of the attention mechanism unaffected. This method helps disperse outlier values because each entry in $\mathbf{Q} \mathbf{M}$ or $\mathbf{K} \mathbf{M}$ is essentially a random weighted sum of the corresponding entries in the original matrices, which lowers quantization-induced errors. In our implementation, inspired by \citet{chee2024quip} and~\citet{tseng2024quip}, the matrix $\mathbf{M}$ is constructed as the product of random diagonal matrices with entries $\pm 1$ and a Hadamard matrix, enabling the multiplication to be computed in $O(d \log d)$ time rather than $O(d^2)$. This step can also be efficiently merged with the rotary embedding computation, thereby avoiding extra computational costs.

\end{enumerate}

We demonstrate empirically that these strategies are effective, lowering numerical error by up to 2.6$\times$, as presented in \cref{sec:numerical_error}.

\section{Empirical Validation}
\label{sec:exp}

To implement \textsc{FlashAttention-3} and assess its performance and precision, we leverage CUTLASS~\citep{Thakkar_CUTLASS_2023} primitives, including the WGMMA and TMA abstractions.

\begin{itemize}

\item \textbf{Benchmarking attention.} We assess the execution time of \textsc{FlashAttention-3} across a range of sequence lengths and juxtapose its performance with a baseline PyTorch implementation, \textsc{FlashAttention-2}, the Triton variant of \textsc{FlashAttention-2} (which is tailored to utilize H100-specific instructions), as well as a vendor-provided \textsc{FlashAttention-2} optimized for H100 GPUs through cuDNN. Our results indicate that \textsc{FlashAttention-3} attains up to a 2.0$\times$ speedup relative to \textsc{FlashAttention-2} and achieves 1.5$\times$ higher performance over the Triton implementation of \textsc{FlashAttention-2}. Furthermore, \textsc{FlashAttention-3} is capable of reaching throughput levels of up to 740 TFLOPs/s, corresponding to 75\% of the H100 GPUs' peak theoretical performance.
\item \textbf{Ablation study.} We demonstrate that our enhancements—specifically, warp-specialization and GEMM-softmax pipelining—are critical contributors to the accelerated runtime of \textsc{FlashAttention-3}.
\item \textbf{Accuracy of FP8 attention.} We show that incorporating block quantization and incoherent processing lowers the numerical inaccuracies of FP8 \textsc{FlashAttention-3} by a factor of 2.6$\times$.
\end{itemize}

\subsection{Benchmarking Attention}
\label{subsec:benchmark_attn}

We evaluate the performance of various attention mechanisms on an H100 80GB SXM5 GPU across multiple configurations (with and without causal masking, head dimensions set to 64 or 128) using FP16 data types. The outcomes, detailed in~\cref{fig:benchmark_attn_fwd} and~\cref{fig:benchmark_attn_bwd}, reveal that \textsc{FlashAttention-3} achieves forward pass speeds roughly 1.5 to 2 times faster than those of \textsc{FlashAttention-2}, and is approximately 1.5 to 1.75 times faster during the backward pass. In comparison to conventional attention implementations, \textsc{FlashAttention-3} attains speedups ranging from 3 up to 16 times. Notably, for sequence lengths of 1k tokens or more, \textsc{FlashAttention-3} also outperforms a proprietary, highly-optimized vendor library (cuDNN) specifically tuned for H100 hardware.

\paragraph{Benchmark settings:} We experiment with sequence lengths of 512, 1k, up to 16k, adjusting the batch size in each case such that the cumulative token count remains at 16k. The hidden dimension is fixed at 2048, while the head dimension is chosen from 64, 128, or 256 (corresponding to 32, 16, or 8 attention heads, respectively). For FLOPs calculation during the forward pass, we employ:

\begin{equation*}
  4 \cdot \text{seqlen}^2 \cdot \text{head dimension} \cdot \text{number of heads}.
\end{equation*}

By applying causal masking, we halve this figure because, on average, only about fifty percent of the values are actually computed. To determine the backward pass FLOPs, we take the FLOPs from the forward pass and scale them by a factor of 2.5—reflecting that the backward step involves five matrix multiplications, whereas the forward step requires only two (with the increase attributable to recomputation).

Additionally, we evaluate the FP8 forward pass runtime using comparable conditions. The outcomes for headdim 256 are shown in ~\cref{fig:benchmark_attn_fp8}, while comprehensive results can be found in \cref{sec:benchmark_attn_fp8_full}.

\subsection{Ablation Study: 2-Stage Pipelining Experiments}
\label{sec:2-stage-pipelining-experiments}

We perform an ablation study on both the two-stage WGMMA-softmax pipeline and warp-specialization using non-causal FP16 \textsc{FlashAttention-3}, holding the parameters constant at $\{\text{batch}, \text{seqlen}, \text{nheads}, \text{hdim}\} = \{ 4, 8448, 16, 128\}$. As shown in~\cref{table:ablation_pipelining}, the results indicate that the introduced algorithmic enhancements—specifically asynchrony through warp-specialization and the overlap of GEMM and softmax—substantially improve performance, increasing throughput from 570 to 661 TFLOPs.

\subsection{Numerical Error Validation}
\label{sec:numerical_error}

Given the growing attention to numerical error in \textsc{FlashAttention}~\citep{golden2024flash}, we conduct a comparative analysis involving \textsc{FlashAttention-2}, \textsc{FlashAttention-3}, and a typical attention mechanism, using a reference FP64 implementation as the baseline. To emulate the presence of outlier features and activations commonly observed in LLMs~\citep{dettmers2208llm, sun2024massive}, we draw the elements of $\mathbf{Q}, \mathbf{K}, \mathbf{V}$ from the following probability distribution:

\begin{equation*}
  \mathcal{N}(0, 1) + \mathcal{N}(0, 100) \cdot \mathrm{Bernoulli}(0.001).
\end{equation*}

Specifically, every element is sampled from a normal distribution with a mean of zero and a standard deviation of 1; however, for 0.1\% of the elements, we incorporate an additional independent term sampled from a normal distribution with a standard deviation of 10. The root mean squared error (RMSE) is subsequently reported in~\cref{table:numerical_error}. When using FP16, both \textsc{FlashAttention-2} and \textsc{FlashAttention-3} yield RMSE values that are 1.7$\times$ lower than those of the conventional implementation, as the intermediate softmax computations are maintained in FP32. In the FP8 baseline attention, per-tensor scaling is applied, the matrix multiplication accumulator operates in FP32, and intermediate softmax values are retained in FP16. Owing to the combination of block quantization and non-coherent computation, \textsc{FlashAttention-3} in FP8 demonstrates a 2.6$\times$ improvement in accuracy relative to the baseline.

\section{Dicussion, Limitations, Conclusion}
\label{sec:discussion}

Through the development of \textsc{FlashAttention-3}, we illustrate how leveraging advanced programming methods and hardware innovations—including asynchrony and reduced-precision computations—can substantially enhance both the speed and reliability of attention mechanisms. Our approach achieves a 1.5-2.0$\times$ acceleration over \textsc{FlashAttention-2}, and it diminishes FP8 numerical errors by a factor of 2.6 when compared to conventional per-tensor quantization approaches. Several avenues remain for future improvement: refining optimization for LLM inference, incorporating a persistent kernel structure within the FP8 kernel,\footnote{For benchmarking purposes, the FP16 \textsc{FlashAttention-3} employs a persistent kernel and an associated load balancing scheme, whereas the FP8 \textsc{FlashAttention-3} does not, which partially accounts for its inferior performance on small sequence lengths and causal masking relative to the FP8 cuDNN kernels.} and gaining a deeper understanding of the role of low-precision attention during large-scale training. While our experimentation centers on Hopper GPUs, we anticipate that the techniques pioneered in this study are broadly applicable to other hardware accelerators as well. Ultimately, we aspire for these advances in attention computation—combining greater efficiency and improved accuracy—to catalyze progress in tasks that involve processing long contexts.

\end{document}